% auto-generated by tools/tablegen_p1.py (from output/validation/kpis.json)
% MANUALLY EDITED -- do not regenerate without re-applying the changes marked CHANGED below.
\begin{table*}[t]
\centering
% CHANGED: caption now states that only two rows are gated MILP outputs, matching the head
% paragraph of Section~\ref{sec:validation}; the previous wording let the reader count four.
% CHANGED: bias/scatter units named in the caption, since the temperature rows are differences.
\caption{Validation against the Memmingen measurements. \emph{Model} is the formulation each
metric is scored against; bias is the signed mean error and scatter the RMSE, both in the unit of
the quantity (K for temperature differences). Two rows are gated outputs of the MILP formulation
that produces the cost results, namely the annual delivered-energy error and the per-step
energy-balance closure; the COP band is a consistency check, the return temperature a
parameter-calibration check, and \emph{n/a} marks the supply-temperature gates, which score the
propagating (NLP) formulation only and cannot be satisfied by a formulation that does not
produce the quantity.}
\label{tab:validation}
\footnotesize
\setlength{\tabcolsep}{4pt}
\begin{tabular}{@{}p{4.9cm}lrrrlc@{}}
\hline
Quantity & Model & Value & Bias & Scatter & Gate & Status \\
\hline
Annual delivered-energy error & MILP & 1.23\,\% & -- & -- & $\le$2\,\% & \checkmark \\
Heat-pump COP (mean) & MILP & 2.99 & -0.00 & -- & $[2.5,5.5]$ & \checkmark \\
% CHANGED: T_return is a fixed MILP parameter (return_temp_c: 63.6), not a solved variable --
% thermal_node.py falls to pyo.Param without return_temp_range/load_factor, and the provenance
% JSON states the simulated series is the nominal constant. The 2.08 K MAE therefore measures
% the spread of the measurement about the calibrated constant, not model accuracy. Reclassified
% from gated validation to parameter-calibration check.
Source return-temp calibration (const.\ param.) & param. & 2.08\,K & $-0.31$\,K & 2.87\,K & -- & report \\
% CHANGED: max split out with its cause named. The 99 % is a vanishing-denominator artefact at
% near-zero flow, not a model failure; the annual closure above is unaffected.
Energy-balance closure (per-step mean) & MILP & 6.4\,\% & -- & -- & $\le$2\,\% & $\times$ \\
\quad per-step max (near-zero flow) & MILP & 99\,\% & -- & -- & -- & report \\
Flow MAPE, source & MILP & 33.8\,\% & -- & -- & -- & report \\
Demand MAPE, total & MILP & 29.1\,\% & -- & -- & -- & report \\
Flow tracking ($r_\text{lvl}$; $r_{\Delta,\text{daily}}$) & MILP & 0.91 / 0.80 & -- & h.\ $\Delta$ 0.24 (noise) & -- & \checkmark \\
% CHANGED: units K on bias/scatter (temperature differences, not levels).
Far-end supply-temp MAE & NLP & 9.21\,K & $+9.21$\,K & 11.31\,K & $\le$1.5\,K & n/a \\
% CHANGED: status n/a and value "--" rather than 0. The baseline does not propagate supply
% temperature, so the trunk drop is zero by construction; the gate was mis-specified and no
% calibration can satisfy it. Reporting a numeric 0 invites reading it as a measured result.
Trunk supply-temp drop MAE & NLP & -- (not propagated) & -- & -- & $\le$1.0\,K & n/a \\
Corridor supply-temp MAE & NLP & 8.32\,K & $-7.73$\,K & -- & -- & n/a \\
\hline
\end{tabular}
\end{table*}